% !TeX root = ../../negocio-tecnologico-exemplo-undf.tex
\chapter{Introdução e oportunidade}
Este relatório apresenta o desenvolvimento e a análise didática de um modelo
de negócio para um serviço digital de apoio à organização de trabalhos
acadêmicos. A proposta foi delimitada à composição de documentos, ao controle
de versões e à orientação de uso da ferramenta. A autoria intelectual e as
decisões de pesquisa permanecem com o estudante.

O objetivo foi construir um modelo coerente, identificar suas hipóteses e
examinar um cenário financeiro simplificado. A modalidade foi tratada como
relatório final de TCC \citeind{manualundf}. O negócio é hipotético: não houve
abertura de empresa, entrevistas, vendas ou operação. Os nomes apresentados
na abertura do documento são fictícios.

\chapter{Fundamentação e método}
A modelagem utilizou os nove componentes do Business Model Canvas como
ferramenta de síntese e discussão \citeind{strategyzer}. A existência de um
quadro preenchido não foi considerada evidência de demanda. Cada componente
foi formulado como hipótese de funcionamento, cuja confirmação exigiria
investigação específica.

O método compreendeu consulta documental à apresentação do LaTeX, definição
do escopo do serviço, elaboração do modelo e análise aritmética de três
cenários. O LaTeX dispõe de mecanismos de composição, referências cruzadas
e bibliografia \citeind{latexproject}; a consulta serviu para delimitar a
base tecnológica. Não foram coletados dados de mercado ou de participantes.

\chapter{Análise do contexto e do mercado}
O público hipotético foi delimitado a estudantes que já escrevem seus
trabalhos e procuram apoio na organização dos arquivos. A necessidade
comercial é uma hipótese deste exercício, sem estimativa de tamanho de
mercado. A disponibilidade gratuita do LaTeX e de modelos acadêmicos foi
considerada uma alternativa relevante à contratação de apoio.

A análise distinguiu disponibilidade técnica de disposição a pagar. A fonte
documental consultada descreve a ferramenta, mas não demonstra que estudantes
contratariam o serviço. Não foram atribuídos percentuais de adesão, preços
observados ou vantagens competitivas comprovadas. O contexto brasileiro
específico requer pesquisa adicional e não foi caracterizado por dados neste exemplo.

\chapter{Solução e base tecnológica}
A solução proposta combina modelos LaTeX, orientações de organização de
arquivos e atendimento remoto de dúvidas de composição. O estágio alcançado
foi a especificação documental do serviço e de suas dependências. Não foi
desenvolvida uma plataforma comercial nem realizado um piloto operacional.

As dependências identificadas incluem distribuição TeX, manutenção dos modelos,
suporte a falhas de compilação e tratamento adequado dos arquivos recebidos.
O serviço foi delimitado para não oferecer escrita de conteúdo acadêmico
em nome do estudante nem prometer aprovação ou conformidade automática
com requisitos de qualquer instituição.

\chapter{Modelo de negócio consolidado}
O quadro a seguir sintetiza o modelo produzido. Todos os componentes são
hipóteses didáticas; nenhum representa comportamento de clientes observado.

\begin{undftabela}[H]{Componentes do modelo de negócio proposto}{tab:modelo-negocio}[Elaboração própria, com organização conceitual baseada em \citeind{strategyzer}]
\begin{tabular}{@{}p{0.29\linewidth}p{0.65\linewidth}@{}}
\hline
\textbf{Componente} & \textbf{Hipótese adotada} \\
\hline
Segmentos & Estudantes que precisam de apoio à composição documental. \\
Proposta de valor & Organização de arquivos e suporte ao uso dos modelos. \\
Canais & Página informativa e atendimento remoto agendado. \\
Relacionamento & Orientação delimitada por sessão e documentação de apoio. \\
Receitas & Cobrança por unidade de atendimento. \\
Recursos & Modelos, ambiente de compilação e conhecimento técnico. \\
Atividades & Suporte, manutenção e atualização das orientações. \\
Parcerias & Colaboradores de suporte, caso a demanda justifique. \\
Custos & Tempo de atendimento e estrutura operacional. \\
\hline
\end{tabular}
\end{undftabela}

A cobrança por atendimento exige delimitação de duração e escopo. A necessidade
de suporte também cria dependência de capacidade de trabalho: aumentar
o número de atendimentos pode elevar o custo e reduzir a disponibilidade.
Assim, receitas, atividades e recursos não foram tratados como escolhas
independentes. O modelo consolidado corresponde à proposta documentada ao
fim do exercício, sem significar validação comercial.

\chapter{Resultados da investigação e avaliação}
A investigação documental confirmou a disponibilidade de uma base técnica
para composição de documentos. Não permitiu concluir que a proposta possui
demanda ou diferenciação suficiente para sustentar um negócio. A análise do
modelo levou à delimitação do serviço ao suporte de composição, evitando
confusão com produção de conteúdo acadêmico por terceiros.

O resultado da avaliação foi um conjunto de hipóteses explícitas e condições
a examinar. Permaneceram sem teste a disposição a pagar, a duração média
de atendimento, o custo de aquisição de clientes e a capacidade operacional.
A ausência desses dados foi registrada como resultado inconclusivo sobre
viabilidade comercial, sem preenchê-la com entrevistas ou métricas fictícias.

\chapter{Viabilidade, riscos e limitações}
Para examinar apenas a consistência aritmética, foram adotadas premissas
\textbf{inteiramente hipotéticas}: preço de R\$~90 por atendimento, custo
variável de R\$~30 por atendimento e custo fixo de R\$~600 por período.
Esses números não são cotações, pesquisa de preços ou estimativa de demanda.
O período é uma unidade de análise comum aos custos e às quantidades.

O resultado simplificado foi calculado por $R=q(p-c)-F$, em que $q$ é a
quantidade de atendimentos, $p$ o preço, $c$ o custo variável unitário e $F$
o custo fixo. A margem unitária é R\$~60 e o equilíbrio aritmético ocorre
em $q=600/60=10$ atendimentos por período.

\begin{undftabela}[H]{Cenários financeiros exclusivamente didáticos}{tab:cenarios-negocio}[Elaboração própria, a partir das premissas hipotéticas declaradas]
\begin{tabular}{@{}rrrr@{}}
\hline
\textbf{Quantidade} & \textbf{Receita (R\$)} & \textbf{Custo total (R\$)} & \textbf{Resultado (R\$)} \\
\hline
5 & 450 & 750 & $-300$ \\
10 & 900 & 900 & 0 \\
15 & 1.350 & 1.050 & 300 \\
\hline
\end{tabular}
\end{undftabela}

A simplificação não detalha tributos, remuneração do trabalho, investimento
inicial ou financiamento. Seu alcance é ilustrativo. O risco de subestimar
tempo e suporte permaneceu aberto, assim como a capacidade de captar dez
atendimentos no período. A viabilidade técnica do serviço completo também
depende de implementação e teste operacional, que não foram realizados.

\chapter{Conclusão}
O objetivo foi atendido pela construção documentada do modelo e pela análise
de suas hipóteses e de um cenário aritmético. O resultado permite compreender
como a proposta funcionaria e quais condições não foram demonstradas.
O equilíbrio calculado não comprova viabilidade comercial.

A análise concluiu que não há evidência suficiente para recomendar a
implantação do negócio. Como continuidade, são pertinentes investigação
com potenciais usuários, teste de escopo do atendimento e levantamento de
custos reais. Essas recomendações decorrem dos limites encontrados; o
relatório registra o trabalho concluído e não um cronograma de execução futura.
